Job advertisements in Europe are spread across government portals, job websites, and company career pages. DataForge is a live data system on Amazon Web Services (AWS). It already collects job data from several sources and stores it in three layers: Bronze (raw data), Silver (clean history), and Gold (reports). It also offers public Jobs and Metrics APIs. Companies now want more than lists of jobs. They want \emph{augmented analytics}: computer methods that add extra labels to jobs and that match a resume to suitable vacancies. If these AI features are added without tests, they can damage data history, increase cost, and produce results that other people cannot repeat.

This thesis designs, builds, and tests a low-cost generative-AI layer on the existing DataForge system. The layer includes: a model gateway with a stop switch and a budget; rule-based data cleaning first; a language-model path that is \textbf{OpenAI-first} in production while Amazon Bedrock quota stays at zero; hybrid search that combines keyword search (BM25) and meaning-based search (embeddings); a Match API that hides personal data and shows evidence from the job text; and an optional small multi-agent process with a fixed limit. Four research questions study: safe integration (RQ1), model and provider choice (RQ2), matching quality versus the old keyword wizard (RQ3), and cost per job (RQ4).

The study uses design science research with one embedded case. The literature review covers lakehouse design, slowly changing dimensions, information retrieval, LLM operations, and European law. On an expanded labelled test set (103 queries and 96 jobs), meaning-based search reaches the highest nDCG@10 (0.2174) in the packaging run. The old keyword wizard scores 0.1462, and BM25 scores 0.1351. Hybrid retrieval scores 0.1668 on this set. A continuous-integration pin to local tf-idf embeddings later records dense nDCG@10 of 0.1258; that pin is reported as a measurement difference, not as a hidden rewrite of the table. A modelled OpenAI-first path for 1{,}000 jobs (rules first, 30\,\% language-model calls, plus embeddings) is about \EUR{0.10}. Live RQ2 on OpenAI \texttt{gpt-4o-mini} completed 40/40 jobs at about \$0.0011 and 1.1\,s average latency. Amazon Nova Micro remains blocked at zero tokens per minute. Nightly enrichment is on for OpenAI (sample rate 0.25), not for Bedrock.

The thesis argues that generative AI in a production lakehouse must be treated as a controlled pipeline stage. History contracts, non-AI baselines, and honest reporting of missing measurements are as important as model choice.

\vspace{0.8em}
\noindent\textbf{Keywords:} lakehouse; augmented analytics; LLM operations; hybrid search; job matching; cost; responsible AI; European job data.
